\section{Codes on time--space diagrams}\label{sec:codes}

The block-level lower bounds of the previous sections (the cut lemmas for unions of slots, Discrete
Isoperimetry and Theorem~M*) use the transition function only through block locality: what they use is
that information flows along the block network of the run, with bounded capacity per edge. This section
makes that structure explicit. A \emph{time--space diagram} records only head motions and an output
schedule; its block network is a directed acyclic graph with bounded degree; and a \emph{causal network
code} on it is an arbitrary assignment of functions to edges. A machine that is correct on a trace class
yields such a code (\cref{lem:class-code}), and conversely these block-level results hold for codes
(\cref{lem:cutset,cor:isoperimetry-code,thm:mstar-diagram}); the cell-level results (the Locality Lemma,
Theorems~P and~P$'$) are not claimed for codes. \Cref{sec:conjecture} then asks whether such codes exist
at all.

\subsection{Diagrams and block networks}

\begin{definition}[Time--space diagram]\label{def:diagram}
\formal{Tape/DefsCE.lean}{TSDiagram, TSDiagram.Valid, runDiagram}
A $k$-tape \emph{time--space diagram} consists of a length $T$, walks $h_t:\N\to\Z$ for $t<k$ (the
work heads), a walk $h_{\mathrm{in}}$ (the input head) and an output schedule $\mathrm{out}:\N\to\N$
($\mathrm{out}(\tau)$ is the number of output bits before step $\tau$). It is \emph{valid} for output
length $L$ if all walks start at $0$ and have unit steps, and the output schedule starts at $0$,
grows by at most one per step, and equals $L$ from time $T$ on. The diagram of a run of $M$ on $w_0$
up to time $T$ consists of its head trajectories and output lengths.
\end{definition}

\begin{definition}[Block network]\label{def:blocknet}
\formal{Tape/DefsCE.lean}{TSDiagram.blockSize, TSDiagram.nodes, TSDiagram.memEdges, TSDiagram.chordEdges, TSDiagram.edges, TSDiagram.rowsAt, TSDiagram.outNode, TSDiagram.visible}
Let $w=b(d)$, $L=d^2w$, $G\ge1$ and $r=\lceil dw/G\rceil$. The block network of a diagram has as
nodes the slots $y<\lceil T/r\rceil$, in their causal order; \emph{memory edges} $y\to y+1$; and a
\emph{chord edge} $x\to z$ whenever some work tape has a chord from $x$ to $z$ of its pair
residency. The rows \emph{attached} at node $y$ (its sources) are the input rows that the input walk
touches during slot $y$; the \emph{sink} of output position $o$ is the slot during which bit $o$ is
appended.
\end{definition}

Every node has in-degree at most $k+1$: one memory edge, and at most one chord per tape by (P2)
(\cref{lem:P2}). There are at most $kN$ chord edges (\lean{card_chordEdges_le}).

\begin{definition}[Causal codes, complete exchange, capacities]\label{def:code}
\formal{Tape/DefsCE.lean}{NetCode, NetCode.inVals, NetCode.edgeVal, NetCode.outVal, TSDiagram.Solves, TSDiagram.CapProfile, TSDiagram.blockCap, TSDiagram.profileCap}
A \emph{causal code} assigns to every edge $e=(u,v)$ a function $\mathrm{enc}_e$ computing a natural
number from the values on the edges into $u$ and the input bits visible at $u$ (the bits of the rows
attached at $u$), and to every output position $o$ a function $\mathrm{dec}_o$ computing a bit from
the same data at the sink of $o$. The code is evaluated along the spine. It \emph{solves complete
exchange} on a set $F$ of inputs within capacities $\mathrm{cap}$ if for every input in $F$ every edge
value is below $2^{\mathrm{cap}(e)}$ and every output bit is the transposed input bit. The capacities
are \emph{within the block profile} if a memory edge $y\to y+1$ carries at most the blocks resident at
slot $y$ (summed over the tapes, at most $3k$) and a chord edge carries one block per tape having that
chord, a block being $rG$ bits. The \emph{profile capacities} are the largest such capacities; with
$s$ extra bits on every memory edge we speak of the block capacities with $s$ extra bits.
\end{definition}

Being a solving code is monotone in the capacities (\lean{solves_mono}), so ruling out codes at the
profile capacities rules them out within every capacity in the profile.

The code may compute arbitrary functions. The only restrictions are causality, the capacities, and
the locality of the inputs: a row is visible only at the nodes where the input walk touches it.

\subsection{The cut-set bound}

\begin{lemma}[Cut-set bound for codes]\label{lem:cutset}
\formal{Tape/Statements/Coding.lean}{cutset_bound_bitwise, cutset_bits_bitwise, cutset_bound, card_le_of_solves}
Let a causal code solve complete exchange on $F$ within $\mathrm{cap}$, and let $U$ be a set of
nodes. Let $B_U$ be the set of input bits whose row is attached at no node of $U$ and whose
transposed position has its sink in $U$. Then
\[
  2^{\abs{B_U}}\,\abs{F}\;\le\;2^{L}\cdot2^{\sum_{e\text{ entering }U}\mathrm{cap}(e)} ,
\]
so if $F$ has co-entropy at most $D$, then $\abs{B_U}\le D+\sum_{e\text{ entering }U}\mathrm{cap}(e)$.
In particular $2^{w\abs{S_U}}\abs{F}\le2^L2^{\sum\mathrm{cap}}$, where $S_U$ is the set of pairs of a row
not attached in $U$ and a column all of whose positions have their sinks in $U$. Taking $U$ to be all
nodes: rows attached nowhere cost co-entropy, $2^{u\,dw}\abs{F}\le2^L$ for $u$ such rows.
\end{lemma}

\begin{outline}
Group the inputs of $F$ by their bits outside $B_U$. Within a group, the values computed inside $U$
are a function of the values on the edges entering $U$ and of the visible bits in $U$, which do not
depend on $B_U$; the outputs in $U$ determine $B_U$. So each group injects into the values on the
entering edges.
\end{outline}

\begin{corollary}[Discrete Isoperimetry for codes]\label{cor:isoperimetry-code}
\formal{Tape/Statements/Coding.lean}{isoperimetry_code, sum_cross_blockCap_le, cutIneq_of_solves}
Let a causal code solve complete exchange on a set of co-entropy at most $D$, on the block network
of a valid diagram with $r\ge1$. For every labelling of the nodes with $\Mb$ label changes and at
most $m$ nodes per label,
\[
  (d\dotminus\Mb)(d\dotminus2m)\,w\;\le\;(\Mb+1)D+\sum_{e:\;\text{labels of the ends differ}}\mathrm{cap}(e),
\]
where $x\dotminus y=\max(x-y,0)$.
At the block capacities with $s$ extra bits, the right-hand sum is at most
$\Mb s+rG(3k\Mb+\chicross)$, so the walks of the diagram satisfy the cut inequality
(\cref{def:cutineq}) with slack $D$ and boundary cost $s$.
\end{corollary}

\begin{outline}
Sum the bits form of \cref{lem:cutset} over the label classes, bound the left side as in
\cref{cor:isoperimetry}, and observe that an edge between different labels is a memory edge at a
label change or a crossing chord.
\end{outline}

\subsection{Distances, flows and the code gap}

\begin{definition}[Distances and routings]\label{def:routing}
\formal{Tape/DefsCE.lean}{ball, netDist, TSDiagram.rowDist, IsWalk, Routing, TSDiagram.RoutesAt}
$\mathrm{dist}$ is the undirected hop distance in the block network. For an input bit $\alpha$,
$\mathrm{rowDist}(\alpha)$ is the distance from the nearest node at which its row is attached to the sink
of its transposed position ($0$ if the row is attached nowhere). A \emph{routing at rate $\lambda$}
within $\mathrm{cap}$ is a fractional multicommodity flow, one commodity per input bit whose row is
attached somewhere, sending $\lambda$ units of each commodity along walks from a node attaching its
row to its sink, with total traversals of every edge at most its capacity.
\end{definition}

\begin{lemma}[Weak duality and volume]\label{lem:volume}
\formal{Tape/Statements/Coding.lean}{weak_duality, rowDist_sum_ge, card_chordEdges_le, sum_cap_le}
A routing at rate $\lambda$ satisfies $\lambda\sum_\alpha\mathrm{rowDist}(\alpha)\le\sum_e\mathrm{cap}(e)$. For a
valid diagram with $1\le r\le dw$ and every $j$,
\[
  j\,L\;\le\;\sum_\alpha\mathrm{rowDist}(\alpha)+j\bigl(4N(2k+3)^{j-1}w+u\,dw\bigr),
\]
where $N$ is the number of nodes and $u$ the number of rows attached nowhere. There are at most $kN$
chord edges, and capacities within the profile sum to at most $4kNrG$.
\end{lemma}

\begin{outline}
Weak duality: each unit of commodity $\alpha$ traverses at least $\mathrm{rowDist}(\alpha)$ edges. For the
distance bound, the ball of radius $j-1$ around a node attaching a row has at most $(2k+3)^{j-1}$ nodes
(degree at most $2k+2$), each node attaches at most two rows, and a sink receives at most $2w$ bits of a
row, so few bits are within distance $j-1$ of their sinks.
\end{outline}

\begin{theorem}[The code gap]\label{thm:code-gap}
\formal{Tape/Statements/Coding.lean}{volume_gap, flow_rate_lt, code_gap}
Fix $k$, $G\ge1$, $c$ and $C$, and let $E=4\clog{(2k+3)}(4k+1)G(c+1)$. For all large $d$, on every valid
$k$-tape diagram of length at most $cL$ with capacities within the profile and at most $C$ rows
attached nowhere,
\[
  w\sum_e\mathrm{cap}(e)\;<\;E\sum_\alpha\mathrm{rowDist}(\alpha),
\]
so every routing has rate $<E/w$. In particular, if a causal code solves complete exchange on a set of
co-entropy at most $Cdw$ within such capacities, the inequality holds: the code (rate~$1$) beats
multicommodity flow by a factor larger than $w/E=\Omega(\log d)$.
\end{theorem}

\begin{outline}
Take $j=\lfloor(w-1)/(2\clog{(2k+3)})\rfloor$, so that $(2k+3)^{2j}\le d-1$. By \cref{lem:volume}, all
but $o(L)$ bits are at distance at least $j=\Omega(w)$ from their sinks, while
$\sum\mathrm{cap}\le4kG(c+1)L$. Weak duality bounds the rate. For a solving code, co-entropy $Cdw$
leaves at most $C$ rows attached nowhere (\cref{lem:cutset}).
\end{outline}

So a code on a block network of a linear-time diagram, if it exists, beats routing by a factor
$\Omega(\log d)$. The connection goes back to Adler, Harvey, Jain, Kleinberg and
Lehman~\cite[Thm.~6.2]{AHJKL2006}, who turn an oblivious algorithm for transposition into a network code on
its computation graph; here \cref{lem:class-code,thm:fast-codes} turn every linear-time transposer into such a
code on its own block network (one way only: a code need not come from a machine). The Tape Coding
Conjecture says that no such code exists.

\subsection{Machines give codes}

\begin{lemma}[Class codes]\label{lem:class-code}
\formal{Tape/Statements/Coding.lean}{traceClass_code}
Let $d\ge1$, let $F$ be a trace class with representative $w_0$ (slot length $r=\lceil dw/G\rceil$,
$G=\max(1,\clog{(g+1)})$) on which $M$ is correct. Then the diagram of the run on $w_0$ up to $T$ is
valid, and its block network carries a causal code that solves complete exchange on $F$, where a
memory edge $y\to y+1$ carries the state, the $k+1$ head positions and the contents of the resident
blocks, at most $\clog{(q+1)}+(k+1)\clog{(2T+1)}+r\clog{(g+1)}\sum_t\abs{R_t(y)}$ bits, a chord edge
carries $r\clog{(g+1)}$ bits per tape having that chord, and no other pair carries anything.
\end{lemma}

\begin{outline}
This is the Block Locality Lemma in functional form. A memory edge $y\to y+1$ carries the state, the head
positions and the contents of the resident blocks at time $(y+1)r$, and a chord edge the contents of its
blocks at the time of its target slot. The code at a node chooses a member of $F$ consistent with the values
it receives and its visible input bits, and sends that member's payloads; the decoder outputs that member's
bit. By \cref{lem:block-locality} for the single slot, all consistent members have the same outgoing payloads
and output bits, and an induction over the slots shows that on every input of $F$ each edge carries that
input's own payload.
\end{outline}

\begin{theorem}[Fast machines give codes]\label{thm:fast-codes}
\formal{Tape/Statements/Coding.lean}{exists_popular_traceClass_on, codes_of_fastOn_bits, solves_profileCap_addTapes, codes_of_fastOn}
Let $M$ be correct within $cL$ steps on every input of a set $F_0$ of co-entropy at most $C_0dw$.
\begin{enumerate}[label=(\alph*)]
\item If $s\ge\clog{(q+1)}+(k+1)\clog{(2cL+1)}$, then there are a valid $k$-tape diagram of length at most
$cL$ (in the proof, the diagram of $M$'s run on the representative of a popular class) and a causal code on
it at the block capacities with $s$ extra bits ($G=\max(1,\clog{(g+1)})$, $r=\lceil dw/G\rceil$), solving
complete exchange on a set of co-entropy at most $(C_0+C_M)dw$, where $C_M$ depends only on $M$ and $c$.
\item Extra bits become tapes: a code at the block capacities with $s\le a\,rG$ extra bits is a code at
the profile capacities of the diagram with $a$ additional stationary tapes.
\item Hence, if $2\clog{(q+1)}\le dw$ and $d$ is large enough, there is such a code at the profile
capacities on a diagram with $k+1$ tapes.
\end{enumerate}
\end{theorem}

\begin{outline}
(a) The popular trace inside $F_0$ (\cref{thm:popular}, which holds within any set of inputs) is a
class of co-entropy at most $(C_0+C_M)dw$; apply \cref{lem:class-code}. (b) A stationary tape changes
no node, edge, source, sink or chord, and its single resident block carries the extra bits on every
memory edge. (c) Combine (a) and (b) with $a=1$.
\end{outline}

\subsection{Theorem M* for code-carrying diagrams}

Theorem~M* holds for codes: its proof uses the machine only through the cut inequality.

\begin{theorem}[Theorem M* for code-carrying diagrams]\label{thm:mstar-diagram}
\formal{Tape/Statements/Coding.lean}{theoremMStar_level_diagram, theoremMStar_diagram, theoremMStar_diagram_profile}
Let $d_i\to\infty$, and let valid $k$-tape diagrams $D_i$ of length at most $cL_i$ carry causal codes
solving complete exchange on sets of co-entropy $m_i$, at the block capacities with $s_i$ extra bits,
where $(m_i+s_i)/L_i\to0$. Then the rescaled work walks of the $D_i$ have a uniformly convergent
subsequence on $[0,cG+1]$, and every subsequential limit is $1$-Lipschitz and satisfies $(\star)$. At
the profile capacities ($s_i=0$) every limit also has total excess at least $1$, total range at most
$k(cG+1)-1$, and its paths are not all monotone.
\end{theorem}

\begin{outline}
By \cref{cor:isoperimetry-code} the walks satisfy the cut inequality with slack $m_i$ and boundary cost
$s_i$; apply Theorem~M* for walks (\cref{thm:mstar-walk}).
\end{outline}

\begin{corollary}[Counterexamples to the conjecture obey $(\star)$]\label{cor:cex-star}
\formal{Tape/Statements/Coding.lean}{mstar_of_robustNoCode_counterexamples, mstar_of_not_robustNoCode}
Every family of counterexamples to the Tape Coding Conjecture (\cref{conj:TCC}): valid diagrams of
length at most $cL$ with codes at the profile capacities on sets of co-entropy $Cdw$, along
$d_i\to\infty$, has rescaled work walks whose subsequential limits are $1$-Lipschitz, satisfy $(\star)$,
have total excess at least $1$ and total range at most $k(cG+1)-1$, and are not all monotone. If the
conjecture fails, such a family exists, with a convergent subsequence.
\end{corollary}

\begin{outline}
\Cref{thm:mstar-diagram} with $m_i=Cd_iw_i=o(L_i)$ and $s_i=0$.
\end{outline}
